% Ad Familiares 4.5 (ad-familiares-04-05)
% Source: https://raw.githubusercontent.com/PerseusDL/canonical-latinLit/master/data/phi0474/phi056/phi0474.phi056.perseus-lat1.xml
% Fetched by scripts/fetch_latin.py

% Dateline (Perseus): Scr. Athenis circ. medio m. Mart. a. 709 (45) .

\ciceroLetterOpener{SERVIVS CICERONI S.}

\ciceroSection{1} postea quam mihi renuntiatum est de obitu Tulliae, filiae tuae, sane quam pro eo ac debui, graviter molesteque tuli communemque eam calamitatem existimavi; qui, si istic adfuissem, neque tibi defuissem coramque meum dolorem tibi declarassem. etsi genus hoc consolationis miserum atque acerbum est, propterea quia, per quos ea confieri debet, propinquos ac familiaris, ii ipsi pari molestia adficiuntur neque sine lacrimis multis id conari possunt, uti magis ipsi videantur aliorum consolatione indigere quam aliis posse suum officium praestare, tamen, quae in praesentia in mentem mihi venerunt, decrevi brevi ad te perscribere, non quo ea te fugere existimem, sed quod forsitan dolore impeditus minus ea perspicias.

\ciceroSection{2} quid est quod tanto opere te commoveat tuus dolor intestinus? cogita , quem ad modum adhuc fortuna nobis cum egerit; ea nobis erepta esse, quae hominibus non minus quam liberi cara esse debent, patriam, honestatem, dignitatem honores omnis. hoc 'uno incommodo addito quid ad dolorem adiungi potuit? aut qui non in illis rebus exercitatus animus callere iam debet atque omnia minoris existimare? an illius vicem, cedo, doles?

\ciceroSection{3} quotiens in eam cogitationem necesse est et tu veneris et nos saepe incidimus, hisce temporibus non pessime cum iis esse actum, quibus sine dolore licitum est mortem cum vita commutare! quid autem fuit quod illam hoc tempore ad vivendum magno opere invitare posset? quae res, quae spes, quod animi solacium? ut cum aliquo adulescente primario coniuncta aetatem gereret? licitum est tibi, credo, pro tua dignitate ex hac iuventute generum deligere, cuius fidei liberos tuos te tuto committere putares. an ut ea liberos ex sese pareret, quos cum florentis videret laetaretur, qui rem a parente traditam per se tenere possent, honores ordinatim petituri essent in re publica, in amicorum negotiis libertate sua usuri? quid horum fuit quod non priusquam datum est ademptum sit? at vero malum est liberos amittere. malum ; nisi hoc peius sit, haec sufferre et perpeti.

\ciceroSection{4} quae res mihi non mediocrem consolationem attulit, volo tibi commemorare, si forte eadem res tibi dolorem minuere possit. ex Asia rediens cum ab Aegina Megaram versus navigarem coepi regiones circumcirca prospicere. post me erat Aegina, ante me Megara, dextra Piraeus, sinistra Corinthus, quae oppida quodam tempore florentissima fuerunt, nunc prostrata et diruta ante oculos iacent. coepi egomet mecum sic cogitare: ' hem ! nos homunculi indignamur, si quis nostrum interiit aut occisus est, quorum vita brevior esse debet, cum uno loco tot oppidum cadavera proiecta iacent? visne ' tu te, Servi, cohibere et meminisse hominem te esse natum?' crede mihi cogitatione ea non mediocriter sum confirmatus. hoc idem, si tibi videtur, fac ante oculos tibi proponas. modo uno tempore tot viri clarissimi interierunt, de imperio populi Romani tanta deminutio facta est, omnes provinciae conquassatae sunt in unius mulierculae animula si iactura facta est, tanto opere commoveris? quae si hoc tempore non diei suum obisset, paucis post annis tamen ei moriendum fuit, quoniam homo nata fuerat.

\ciceroSection{5} etiam tu ab hisce rebus animum ac cogitationem tuam avoca atque ea potius reminiscere, quae digna tua persona sunt, illam, quam diu ei opus fuerit, vixisse, una cum re publica fuisse, te, patrem suum, praetorem, consulem, augurem vidisse, adulescentibus primariis nuptam fuisse, omnibus bonis prope perfunctam esse, cum res publica occideret vita excessisse. quid est quod tu aut illa cum fortuna hoc nomine queri possitis? denique noli te oblivisci Ciceronem esse et eum, qui aliis consueris praecipere et dare consilium, neque imitare malos medicos, qui in alienis morbis profitentur tenere se medicinae scientiam, ipsi se curare non possunt, sed potius, quae aliis tute praecipere soles, ea tute tibi subiace atque apud animum propone.

\ciceroSection{6} nullus dolor est, quem non longinquitas temporis minuat ac molliat. hoc te exspectare tempus tibi turpe est ac non ei rei sapientia tua te occurrere. quod si qui etiam inferis sensus est, qui illius in te amor fuit pietasque in omnis suos, hoc certe illa te facere non vult. da hoc illi mortuae, da ceteris amicis ac familiaribus, qui tuo dolore maerent, da patriae, ut, si qua in re opus sit, opera et consilio tuo uti possit. denique , quoniam in eam fortunam devenimus, ut etiam huic rei nobis serviendum sit, noli committere ut quisquam te putet non tam filiam quam rei publicae tempora et aliorum victoriam lugere. plura me ad te de hac re scribere pudet, ne videar prudentiae tuae diffidere. qua re, si hoc unum proposuero, finem faciam scribendi: vidimus aliquotiens secundam pulcherrime te ferre fortunam magnamque ex ea re te laudem apisci; fac aliquando intellegamus adversam quoque te aeque ferre posse neque id maius, quam debeat, tibi onus videri, ne ex omnibus virtutibus haec una tibi videatur deesse. quod ad me attinet, cum te tranquilliorem animo esse cognoro, de iis rebus, quae hic geruntur, quemadmodumque se provincia habeat, certiorem faciam. vale .
